\section{Problem formulation and Galerkin reduction}
\label{sec:problem}

\subsection{Parameterized dynamics and POD coordinates}
\label{sec:pod}

Let $\mu\in\Pset\subset\R^{d_\mu}$ denote a physical parameter. The resolved
velocity and algebraic pressure are represented by
\begin{align}
 u_r(x,t;\mu)
 &=\overline u(x)+\sum_{i=1}^{r}a_i(t;\mu)\varphi_i(x),
 \label{eq:velocity-pod}\\
 p_{r_p}(x,t;\mu)
 &=\overline p(x)+\sum_{\ell=1}^{r_p}b_\ell(t;\mu)\psi_\ell(x),
 \label{eq:pressure-pod}
\end{align}
where $a(t;\mu)\in\R^r$ and $b(t;\mu)\in\R^{r_p}$. Projection of the
momentum equation gives a semi-explicit velocity equation
\begin{equation}
 \dot a=f_G(a,b;\mu)\in\R^r.
 \label{eq:semiexplicit-galerkin}
\end{equation}
The maps and matrices introduced below are fixed analytic or operator-derived
objects. They are not neural-network weights.

\subsection{Semi-explicit pressure treatment}
\label{sec:pressure}

Pressure is recovered from an algebraic Pressure--Poisson map
\begin{equation}
 b=\mathcal P_p(a;\mu).
 \label{eq:pressure-map}
\end{equation}
The pressure-eliminated Galerkin field is
\begin{equation}
 F_G(a;\mu)
 :=f_G\!\left(a,\mathcal P_p(a;\mu);\mu\right).
 \label{eq:pressure-eliminated}
\end{equation}
Thus pressure has no independent time derivative in the present theory.
Equation~\eqref{eq:pressure-map} may be evaluated analytically inside
$F_G$, or after each accepted velocity macro-step in a semi-explicit
implementation.

\subsection{Truncation-induced non-Markovian closure}
\label{sec:nonmarkovian}

Elimination of discarded Galerkin modes generally produces an
initial-condition term and a history operator. We therefore write the
resolved equation as
\begin{equation}
 \dot a(t)=F_G(a(t);\mu)+r_{\mem}(t),
 \label{eq:resolved-closure}
\end{equation}
and seek a finite-dimensional internal realization of $r_{\mem}$. The
purpose of the memory state is to replace an explicit finite delay window,
not to replace the resolved Galerkin field.

\subsection{Assumptions}
\label{sec:assumptions}

The assumptions are stated separately so that each result can invoke only
what it needs.

\paragraph{A1 (compact parameter domain).}
$\Pset\subset\R^{d_\mu}$ is compact.

\paragraph{A2 (Galerkin regularity).}
For each $\mu\in\Pset$, $F_G(\cdot;\mu)$ is locally Lipschitz in $a$,
uniformly with respect to $\mu$ on compact subsets of $\R^r\times\Pset$.

\paragraph{A3 (memory-coefficient regularity).}
The maps $k,v,q,\Gamma,\eta$ are continuous in $(a,\mu)$ and locally
Lipschitz in $a$, uniformly in $\mu$ on compact subsets. The matrix
$\Gamma(a,\mu)\in\R^{d_k\times d_k}$ is real and diagonal.

\paragraph{A4 (normalized addresses).}
For every admissible $(a,\mu)$,
\begin{equation}
 \norm{k(a;\mu)}\le1,\qquad \norm{q(a;\mu)}\le1.
 \label{eq:address-bounds}
\end{equation}

\paragraph{A5 (strict forgetting).}
There is a constant $\gamma_{\min}>0$ such that
\begin{equation}
 \Gamma(a,\mu)\succeq\gamma_{\min}I_{d_k}.
 \label{eq:gamma-lower}
\end{equation}

\paragraph{A6 (bounded write rate).}
There is an $\eta_{\max}<\infty$ such that
\begin{equation}
 0\le\eta(a,\mu)\le\eta_{\max}.
 \label{eq:eta-bound}
\end{equation}

\paragraph{A7 (bounded write target along the trajectory).}
For the memory boundedness result, and only for that result, there is a
$v_{\max}<\infty$ such that
\begin{equation}
 \norm{v(a(t);\mu)}\le v_{\max}.
 \label{eq:v-bound}
\end{equation}
This condition must follow from the construction of $v$ or from a separate
bound on $a(t)$; it is not inferred from A1--A6.

\paragraph{A8 (stable unresolved linear block).}
In the linear resolved--unresolved correspondence,
$A_{uu}$ is Hurwitz.

\paragraph{A9 (reciprocal energy coupling).}
The energy-stable subclass uses
\begin{equation}
 \dot a=F_G(a)-\sum_{j=1}^{m}C_j^\top y_j,\qquad
 \dot y_j=\eta_jC_ja-\lambda_jy_j,
 \quad \eta_j>0,\quad\lambda_j>0.
 \label{eq:reciprocal-model-assumption}
\end{equation}

\section{Continuous delta-memory closure}
\label{sec:continuous-memory}

\subsection{Discrete channel-wise delta recurrence}
\label{sec:discrete-recurrence}

Let $S_n\in\R^{d_k\times d_v}$ be a fixed-capacity matrix state,
$k_n\in\R^{d_k}$ a write address, and $v_n\in\R^{d_v}$ a write target.
The discrete recurrence is
\begin{equation}
 S_{n+1}
 =(I_{d_k}-\beta_nk_nk_n^\top)D_nS_n+\beta_nk_nv_n^\top.
 \label{eq:discrete-delta}
\end{equation}
With $\overline S_n=D_nS_n$, the same update is
\begin{equation}
 S_{n+1}
 =\overline S_n+\beta_nk_n
   (v_n-\overline S_n^\top k_n)^\top.
 \label{eq:erase-write}
\end{equation}
This form separates channel-wise forgetting from reconstruction-error
writing.

\subsection{Physical-time scaling and continuous limit}
\label{sec:continuous-limit}

Let $h_n=t_{n+1}-t_n>0$ and set
\begin{equation}
 D_n=\exp(-h_n\Gamma_n),\qquad
 \beta_n=1-\exp(-h_n\eta_n).
 \label{eq:physical-scaling}
\end{equation}

\phantomsection\label{lem:expansions}
\begin{lemmaone}
If $\Gamma_n$ and $\eta_n$ remain bounded as $h_n\to0$, then
\begin{equation}
 D_n=I_{d_k}-h_n\Gamma_n+O(h_n^2),\qquad
 \beta_n=h_n\eta_n+O(h_n^2).
 \label{eq:first-order-expansions}
\end{equation}
The remainder constants are uniform when $\Gamma_n$ and $\eta_n$ range over
fixed compact sets.
\end{lemmaone}

\paragraph{Proof idea.}
Apply Taylor's theorem to the matrix exponential and the scalar exponential.
The complete estimate is in Appendix~\ref{app:expansions}.

Substitution of \eqref{eq:first-order-expansions} into
\eqref{eq:erase-write} gives
\begin{equation}
 S_{n+1}
 =S_n+h_n\left[
 -\Gamma_nS_n+\eta_nk_n(v_n-S_n^\top k_n)^\top
 \right]+O(h_n^2).
 \label{eq:discrete-first-order}
\end{equation}
The associated continuous memory equation is
\begin{equation}
 \boxed{
 \dot S=-\Gamma S+\eta k(v-S^\top k)^\top.}
 \label{eq:continuous-memory}
\end{equation}

\subsection{CDM-GROM augmented system}
\label{sec:augmented-system}

Let
\begin{equation}
 \phi(a;\mu)=
 [1,\ a^\top,\ (a\otimes a)^\top,\ \mu^\top]^\top
 \in\R^{d_\phi}.
 \label{eq:analytic-features}
\end{equation}
For fixed compatible matrices $K,V,Q$, define
\begin{align}
 k(a;\mu)
 &=\frac{K\phi(a;\mu)}
 {\norm{K\phi(a;\mu)}+\varepsilon}\in\R^{d_k},
 \label{eq:k-map}\\
 v(a;\mu)&=V\phi(a;\mu)\in\R^{d_v},
 \label{eq:v-map}\\
 q(a;\mu)
 &=\frac{Q\phi(a;\mu)}
 {\norm{Q\phi(a;\mu)}+\varepsilon}\in\R^{d_k},
 \label{eq:q-map}
\end{align}
where $\varepsilon>0$ is fixed. With a fixed
$B\in\R^{r\times d_v}$, the memory read and closure are
\begin{equation}
 y=S^\top q\in\R^{d_v},\qquad r_{\mem}=By\in\R^r.
 \label{eq:memory-read}
\end{equation}
The CDM-GROM is
\begin{equation}
 \boxed{
 \begin{aligned}
 \dot a
 &=F_G(a;\mu)+BS^\top q(a;\mu),\\
 \dot S
 &=-\Gamma(a,\mu)S+
 \eta(a,\mu)k(a,\mu)
 [v(a,\mu)-S^\top k(a,\mu)]^\top.
 \end{aligned}}
 \label{eq:cdm-grom}
\end{equation}
The pair $(a,S)\in\R^r\times\R^{d_k\times d_v}$ is a finite-dimensional
Markov state. Elimination of $S$ produces a history-dependent equation for
$a$. Pressure is recovered through \eqref{eq:pressure-map}.

\subsection{Gradient-flow interpretation}
\label{sec:gradient-flow}

\phantomsection\label{prop:gradient-flow}
\begin{propositionone}
Fix $k\in\R^{d_k}$, $v\in\R^{d_v}$, a symmetric positive-semidefinite
$\Gamma\in\R^{d_k\times d_k}$, and $\eta>0$. Define
\begin{equation}
 \mathcal J(S)
 =\frac12\norm{v-S^\top k}^2+
 \frac{1}{2\eta}\tr(S^\top\Gamma S).
 \label{eq:reconstruction-objective}
\end{equation}
Then \eqref{eq:continuous-memory} is the Frobenius gradient flow
\begin{equation}
 \dot S=-\eta\nabla_S\mathcal J(S).
 \label{eq:gradient-flow}
\end{equation}
\end{propositionone}

\paragraph{Proof idea.}
Differentiate both quadratic terms with respect to $S$. Full algebra is in
Appendix~\ref{app:gradient-flow}.

\phantomsection\label{cor:fixed-equilibrium}
\begin{corollarytwo}
Under the conditions of \hyperref[prop:gradient-flow]{Proposition 1}, if
$\Gamma\succ0$, the fixed-input memory equation has the unique equilibrium
\begin{equation}
 S_\star=(\Gamma+\eta kk^\top)^{-1}\eta kv^\top.
 \label{eq:memory-equilibrium}
\end{equation}
Moreover,
\begin{equation}
 \fnorm{S(t)-S_\star}
 \le
 \exp[-\lambda_{\min}(\Gamma+\eta kk^\top)t]
 \fnorm{S(0)-S_\star}.
 \label{eq:equilibrium-convergence}
\end{equation}
\end{corollarytwo}

\section{Finite-state and memory-kernel representations}
\label{sec:finite-state}

\subsection{Fixed-key minimal realization}
\label{sec:fixed-key}

\phantomsection\label{prop:fixed-key}
\begin{propositionthree}
Let $k=q=e\in\R^{d_k}$ with $\norm e=1$, let
$\Gamma=\gamma I_{d_k}$ with $\gamma>0$, let $\eta\ge0$ be constant, and
let $v(a)=Ca$ for $C\in\R^{d_v\times r}$. Define
$y=S^\top e\in\R^{d_v}$. Then
\begin{equation}
 \dot y=-(\gamma+\eta)y+\eta Ca.
 \label{eq:fixed-key-vector}
\end{equation}
If $P_\perp=I_{d_k}-ee^\top$, then
\begin{equation}
 P_\perp S(t)=\exp(-\gamma t)P_\perp S(0).
 \label{eq:orthogonal-decay}
\end{equation}
Consequently, $P_\perp S$ neither affects the read $S^\top e$ nor the
closure $BS^\top e$, and the vector $y$ is a minimal observable state for
the single fixed address, up to ordinary controllability/observability
reductions of the pair $(B,C)$.
\end{propositionthree}

\paragraph{Proof idea.}
Left-multiply the matrix equation by $e^\top$ and by $P_\perp$,
respectively. See Appendix~\ref{app:fixed-key}.

\subsection{Volterra-kernel equivalence}
\label{sec:volterra}

\phantomsection\label{thm:volterra}
\begin{theoremfour}
For $j=1,\ldots,m$, let $y_j(t)\in\R^{d_j}$ satisfy
\begin{equation}
 \dot y_j=-\Lambda_jy_j+C_ja,
 \qquad
 B_j\in\R^{r\times d_j},\quad
 C_j\in\R^{d_j\times r},
 \label{eq:auxiliary-system}
\end{equation}
where every $\Lambda_j$ is exponentially stable in the convention
$\dot y_j=-\Lambda_jy_j$. Then
\begin{equation}
 \dot a(t)=F_G(a(t);\mu)+g_m(t)
 +\int_0^tK_m(t-s)a(s)\dd s,
 \label{eq:volterra-resolved}
\end{equation}
where
\begin{align}
 g_m(t)&=\sum_{j=1}^{m}B_j\e^{-\Lambda_jt}y_j(0),
 \label{eq:initial-slip}\\
 K_m(\tau)&=\sum_{j=1}^{m}B_j\e^{-\Lambda_j\tau}C_j.
 \label{eq:finite-kernel}
\end{align}
For $\operatorname{Re}s$ sufficiently large,
\begin{equation}
 \widehat K_m(s)
 =\sum_{j=1}^{m}B_j(sI_{d_j}+\Lambda_j)^{-1}C_j.
 \label{eq:kernel-laplace}
\end{equation}
\end{theoremfour}

\paragraph{Proof idea.}
Use variation of constants for every $y_j$, substitute into the resolved
equation, and exchange a finite sum with the integral. See
Appendix~\ref{app:volterra}.

\subsection{Linear resolved--unresolved Galerkin correspondence}
\label{sec:linear-correspondence}

\phantomsection\label{thm:linear-elimination}
\begin{theoremfive}
Consider the linear or linearized block system
\begin{equation}
 \begin{bmatrix}\dot a\\\dot c\end{bmatrix}
 =
 \begin{bmatrix}
 A_{rr}&A_{ru}\\
 A_{ur}&A_{uu}
 \end{bmatrix}
 \begin{bmatrix}a\\c\end{bmatrix},
 \label{eq:block-linear}
\end{equation}
where $a\in\R^r$ and $c\in\R^{r_u}$. Eliminating $c$ gives
\begin{equation}
 \dot a(t)=A_{rr}a(t)+A_{ru}\e^{A_{uu}t}c(0)
 +\int_0^tK_{\rm exact}(t-s)a(s)\dd s,
 \label{eq:exact-linear-memory}
\end{equation}
with
\begin{equation}
 K_{\rm exact}(\tau)=A_{ru}\e^{A_{uu}\tau}A_{ur}.
 \label{eq:exact-kernel}
\end{equation}
Assumption A8 is required for exponential decay of the slip and kernel, but
not for the finite-time elimination identity itself.
\end{theoremfive}

\paragraph{Proof idea.}
Solve the unresolved linear equation by variation of constants. The complete
derivation is in Appendix~\ref{app:linear-elimination}.

\phantomsection\label{cor:exact-realization}
\begin{corollarythree}
Under A8, the choice
\begin{equation}
 y=c,\qquad B=A_{ru},\qquad C=A_{ur},\qquad \Lambda=-A_{uu}
 \label{eq:exact-realization}
\end{equation}
is an exponentially stable finite realization that exactly recovers
\eqref{eq:exact-kernel}. A reduced stable realization replaces
$K_{\rm exact}$ by
$K_m(\tau)=\sum_jB_j\e^{-\Lambda_j\tau}C_j$; its resolved-state error is
controlled by \hyperref[thm:kernel-error]{Theorem 7}.
\end{corollarythree}

\subsection{Operator-based construction}
\label{sec:operator-construction}

Three non-neural constructions are available.

\paragraph{Exact unresolved realization.}
Construct an enlarged POD--Galerkin model, partition its linearization into
\eqref{eq:block-linear}, and use \eqref{eq:exact-realization}. This preserves
the exact linear kernel but retains all unresolved linear states.

\paragraph{Reduced unresolved realization.}
Apply spectral truncation, balanced truncation, rational Krylov reduction, or
a stable projection to the unresolved subsystem. Stability of the reduced
$-\Lambda_j$ blocks must be verified rather than inferred.

\paragraph{Stable pole--residue approximation.}
Approximate
\begin{equation}
 G_{\mem}(s)=A_{ru}(sI-A_{uu})^{-1}A_{ur}
 \label{eq:memory-transfer}
\end{equation}
by \eqref{eq:kernel-laplace}, with every pole in the open left half-plane.
Least-squares or vector-fitting estimation is operator/data-assisted
identification, not neural learning.

\section{Mathematical analysis}
\label{sec:analysis}

\subsection{Well-posedness and boundedness}
\label{sec:wellposedness}

\phantomsection\label{thm:local-wellposedness}
\begin{theoremtwo}
Under A1--A3, for every $\mu\in\Pset$ and every initial condition
$(a_0,S_0)\in\R^r\times\R^{d_k\times d_v}$, the CDM-GROM
\eqref{eq:cdm-grom} has a unique maximal local solution. The theorem makes no
unconditional claim of global existence.
\end{theoremtwo}

\paragraph{Proof idea.}
The augmented right-hand side is locally Lipschitz in $(a,S)$ on bounded
sets. Picard--Lindel\"of then applies. See Appendix~\ref{app:wellposedness}.

\phantomsection\label{thm:memory-boundedness}
\begin{theoremthree}
Assume A5--A7 along a solution of \eqref{eq:cdm-grom}. Then
\begin{equation}
 \fnorm{S(t)}^2
 \le
 \e^{-2\gamma_{\min}t}\fnorm{S(0)}^2+
 \frac{\eta_{\max}v_{\max}^2}{4\gamma_{\min}}
 \left(1-\e^{-2\gamma_{\min}t}\right).
 \label{eq:memory-explicit-bound}
\end{equation}
The estimate holds throughout the maximal interval of existence. If the
coupled solution is global, then
\begin{equation}
 \limsup_{t\to\infty}\fnorm{S(t)}^2
 \le\frac{\eta_{\max}v_{\max}^2}{4\gamma_{\min}}.
 \label{eq:memory-limsup}
\end{equation}
\end{theoremthree}

\paragraph{Proof idea.}
With $y=S^\top k$, take the Frobenius product of the memory equation with
$S$ and complete the square in $y$. See Appendix~\ref{app:memory-bound}.

\phantomsection\label{prop:memory-passivity}
\begin{propositiontwo}
Under A5--A6, the exact identity
\begin{equation}
 \frac12\frac{\dd}{\dd t}\fnorm S^2
 =-\ipF{S}{\Gamma S}+\eta y^\top(v-y),
 \qquad y=S^\top k,
 \label{eq:memory-energy-identity}
\end{equation}
implies
\begin{equation}
 \frac12\frac{\dd}{\dd t}\fnorm S^2
 \le
 -\gamma_{\min}\fnorm S^2
 -\frac{\eta}{2}\norm y^2
 +\frac{\eta}{2}\norm v^2.
 \label{eq:passivity-bound}
\end{equation}
\end{propositiontwo}

\begin{remark}
With storage $\fnorm S^2/2$, input $v$, and read output $y=S^\top k$,
\eqref{eq:passivity-bound} is an output-strict dissipation inequality with
supply rate $\eta\norm v^2/2$. It is also an input-to-state estimate when
$v$ is bounded.
\end{remark}

\begin{remark}
\hyperref[thm:memory-boundedness]{Theorem 3} bounds only the memory state along
a trajectory for which A7 holds. It neither bounds $a$ nor prevents a
finite-time escape of the coupled CDM-GROM.
\end{remark}

\subsection{Discrete-to-continuous consistency}
\label{sec:consistency}

\phantomsection\label{thm:local-consistency}
\begin{theoremone}
Assume A3 and suppose the sampled coefficient functions and the exact memory
solution are sufficiently smooth on $[0,T]$ to admit uniform second-order
Taylor remainders. If the exact value $S(t_n)$ is inserted into
\eqref{eq:discrete-delta}, then its one-step defect relative to
\eqref{eq:continuous-memory} satisfies
\begin{equation}
 \fnorm{S_{n+1}^{\rm rec}-S(t_{n+1})}
 \le C_{\rm loc}h_n^2
 \label{eq:local-defect}
\end{equation}
for all sufficiently small $h_n$, with $C_{\rm loc}$ independent of $n$.
\end{theoremone}

\paragraph{Proof idea.}
Combine \hyperref[lem:expansions]{Lemma 1} with a uniform product-remainder
estimate. See Appendix~\ref{app:local-consistency}.

\phantomsection\label{cor:global-first-order}
\begin{corollaryone}
In addition to the conditions of \hyperref[thm:local-consistency]{Theorem 1},
assume that the exact and discrete trajectories remain in a common compact
set on $[0,T]$ and that the one-step map is Lipschitz there with factor
$1+L_hh_n$. Let $h=\max_nh_n$. Then, for identical initial data,
\begin{equation}
 \max_{t_n\le T}\fnorm{S_n-S(t_n)}
 \le
 \frac{C_{\rm loc}}{L_h}
 \left(\e^{L_hT}-1\right)h
 \label{eq:global-first-order}
\end{equation}
when $L_h>0$; the right-hand side is $C_{\rm loc}Th$ when $L_h=0$.
\end{corollaryone}

\subsection{Energy-stable reciprocal subclass}
\label{sec:energy-subclass}

\phantomsection\label{thm:reciprocal-energy}
\begin{theoremsix}
Under A9, define
\begin{equation}
 \mathcal E(a,y_1,\ldots,y_m)
 =\frac12\norm a^2+
 \sum_{j=1}^{m}\frac{1}{2\eta_j}\norm{y_j}^2.
 \label{eq:augmented-energy}
\end{equation}
Along every classical solution,
\begin{equation}
 \frac{\dd\mathcal E}{\dd t}
 =a^\top F_G(a)
 -\sum_{j=1}^{m}\frac{\lambda_j}{\eta_j}\norm{y_j}^2.
 \label{eq:reciprocal-energy-identity}
\end{equation}
Hence $\mathcal E$ is non-increasing on the interval of existence if
$a^\top F_G(a)\le0$. If A2 also holds and, more generally,
\begin{equation}
 a^\top F_G(a)\le c_0-c_1\norm a^2,
 \qquad c_0\ge0,\quad c_1>0,
 \label{eq:galerkin-absorbing}
\end{equation}
then, with
\begin{equation}
 \kappa=2\min\{c_1,\lambda_1,\ldots,\lambda_m\},
 \label{eq:kappa}
\end{equation}
\begin{equation}
 \mathcal E(t)
 \le
 \e^{-\kappa t}\mathcal E(0)
 +\frac{c_0}{\kappa}(1-\e^{-\kappa t}).
 \label{eq:augmented-absorbing-bound}
\end{equation}
The energy bound prevents finite-time escape; consequently the local
solution is global and \eqref{eq:augmented-absorbing-bound} is an absorbing
estimate.
\end{theoremsix}

\paragraph{Proof idea.}
Differentiate \eqref{eq:augmented-energy}; reciprocal cross terms cancel
exactly. See Appendix~\ref{app:reciprocal-energy}.

\subsection{Kernel approximation error}
\label{sec:kernel-error}

\phantomsection\label{thm:kernel-error}
\begin{theoremseven}
Let $a$ and $a_m$ solve, respectively,
\begin{align}
 \dot a(t)
 &=F_G(a(t))+g(t)+\int_0^tK(t-s)a(s)\dd s,
 \label{eq:exact-volterra}\\
 \dot a_m(t)
 &=F_G(a_m(t))+g_m(t)+\int_0^tK_m(t-s)a_m(s)\dd s.
 \label{eq:approx-volterra}
\end{align}
Assume A2 on a common ball containing both trajectories, with Lipschitz
constant $L$, and suppose
\begin{equation}
 \sup_{0\le t\le T}\norm{a_m(t)}\le M,\qquad
 K,K_m\in L^1(0,T),\qquad g-g_m\in L^1(0,T).
 \label{eq:kernel-error-assumptions}
\end{equation}
Define
\begin{align}
 \delta_K(T)&=\int_0^T\norm{K(\tau)-K_m(\tau)}\dd\tau,
 \label{eq:delta-k}\\
 \delta_g(T)&=\int_0^T\norm{g(t)-g_m(t)}\dd t,
 \label{eq:delta-g}\\
 \kappa_T&=\int_0^T\norm{K(\tau)}\dd\tau.
 \label{eq:kappa-t}
\end{align}
Then
\begin{equation}
 \sup_{0\le t\le T}\norm{a(t)-a_m(t)}
 \le
 \left[
 \norm{a(0)-a_m(0)}+\delta_g(T)+MT\delta_K(T)
 \right]\e^{(L+\kappa_T)T}.
 \label{eq:kernel-error-bound}
\end{equation}
Thus one may take
$C_T=\e^{(L+\kappa_T)T}\max\{1,MT\}$ when the three error contributions
are grouped additively.
\end{theoremseven}

\paragraph{Proof idea.}
Subtract the integral forms, use Fubini's theorem for the double convolution
integral, and apply Gronwall's inequality. See
Appendix~\ref{app:kernel-error-proof}.

\subsection{Scope and limitations}
\label{sec:limitations}

The following restrictions are part of the method statement.
\begin{enumerate}
 \item Positive diagonal $\Gamma$ primarily generates decaying time scales.
 Purely real exponential kernels can inefficiently represent strongly
 oscillatory memory. Two-by-two rotational blocks are a possible extension,
 but are not included here.
 \item State-dependent $k$ and $q$ define a nonlinear fading-memory operator;
 the simple convolution formula \eqref{eq:finite-kernel} then does not apply.
 \item The matrix $S$ need not be a minimal realization for fixed addresses,
 as shown by \hyperref[prop:fixed-key]{Proposition 3}.
 \item Memory boundedness is not full-ROM stability.
 \item The correspondence in
 \hyperref[thm:linear-elimination]{Theorem 5} is exact only for a linear or
 linearized system.
 \item The reciprocal subclass can over-constrain energy-growth mechanisms
 near a Hopf bifurcation.
 \item For a nonnormal $\Lambda_j$, spectral stability implies exponential
 decay but not monotone Euclidean energy. A Lyapunov weight is required for a
 direct quadratic decay estimate.
 \item The present work does not claim the first use of memory in ROMs.
 Relevant prior-work citations remain to be supplied:
 \textbf{[TODO citation: finite-memory and Mori--Zwanzig ROM literature]}.
\end{enumerate}
